% generated by experiments/make_table2.py; do not edit by hand
\begin{table}[htbp]
\centering
\caption{真实激光下最好的单环（加大下载端口耦合组，Mackey--Glass，现实噪声的其余各项不变）。失谐抖动由激光的白频率噪声按每个符号平均得到；“固定工作点”是 1\% 抖动假设下选出的点（$Q=1.9\times10^{6}$），“重选”是在该激光的噪声下按第~\ref{sec:protocol}~节的流程重新选点（每条序列 56 个按 $Q$ 分档的候选）。3 条序列 $\times$ 6 对新运行。两列用了不同的噪声实现，同一个工作点在两列之间约有 9\% 的差别。}
\label{tab:laser}
\small
\begin{tabular}{p{4.2cm}cccc}
\toprule
激光本征线宽 & 每符号抖动 & 固定工作点 & 重选后 & 重选点的 $Q$ \\
\midrule
100~Hz & 0.25\% & $6.7\times10^{-5}$ & -- & -- \\
1.6~kHz & 1.00\% & $1.3\times10^{-4}$ & $1.2\times10^{-4}$ & $1.9\times10^{6}$ \\
16~kHz & 3.15\% & $3.1\times10^{-4}$ & $2.9\times10^{-4}$ & $1.9\times10^{6}$ \\
100~kHz & 7.88\% & $8.1\times10^{-4}$ & $6.4\times10^{-4}$ & $1.9\times10^{6}$、$2.5\times10^{6}$ \\
1~MHz & 24.92\% & $5.2\times10^{-3}$ & $1.6\times10^{-3}$ & $1.2\times10^{6}$ \\
\midrule
100~kHz，加闪烁噪声（转折 1~MHz）和 1~MHz 锁定 & 7.92\% & $8.3\times10^{-4}$ & -- & -- \\
1.6~kHz，加闪烁噪声（转折 10~MHz），不锁定 & 1.10\% & $1.8\times10^{-4}$ & -- & -- \\
\bottomrule
\end{tabular}
\end{table}
